% --- Part 2: Fundamentals of HE & Modern Schemes (Slides 4-8) ---
\section{Fundamentos de HE}
\SectionFrame{Fundamentos de HE}

\begin{frame}{Orígenes y Evolución}
    \begin{itemize}
        \item \textbf{1978:} Rivest, Adleman y Dertouzos proponen \textit{"Privacy Homomorphisms"}.
        \[ E(m_1 \otimes m_2) = E(m_1) \otimes E(m_2) \]
        \item \textbf{Cifrado Parcialmente Homomórfico (PHE):}
        \begin{itemize}
            \item \textbf{RSA (Multiplicativo):}
            \[ c_1 \cdot c_2 = m_1^e \cdot m_2^e = (m_1 \cdot m_2)^e \pmod n \]
            \item \textbf{Paillier (Aditivo):}
            \[ g^{m_1} r_1^n \cdot g^{m_2} r_2^n = g^{m_1+m_2} (r_1 r_2)^n \pmod{n^2} \]
        \end{itemize}
    \end{itemize}
\end{frame}

\begin{frame}{Hacia el Cifrado Total (FHE)}
    \begin{itemize}
        \item \textbf{2009: Craig Gentry} presenta el primer esquema Totalmente Homomórfico (FHE).
        \item \textbf{Leveled FHE:} Esquemas diseñados para evaluar circuitos de una profundidad máxima predefinida.
        \item \textbf{Bootstrapping:} La gran innovación de Gentry.
        \begin{itemize}
            \item Proceso de "refresco" del texto cifrado para reducir el ruido.
            \item Evaluación homomórfica del propio circuito de descifrado.
        \end{itemize}
    \end{itemize}
    \vfill
    \begin{figure}
        \centering
        \includegraphics[width=0.6\textwidth]{bootstrap.png}
        \caption{Diagrama esquemático del proceso de Bootstrapping.}
    \end{figure}
\end{frame}

\begin{frame}{El Problema del "Ruido"}
    \begin{itemize}
        \item El cifrado incluye un "error" o ruido para garantizar seguridad.
        \item El ruido crece con cada operación homomórfica.
        \item La multiplicación aumenta el ruido mucho más rápido que la suma.
        \item Si el ruido supera un umbral, el descifrado falla.
    \end{itemize}
    \begin{center}
        \includegraphics[width=0.6\textwidth]{noise.png}
    \end{center}
\end{frame}

\begin{frame}{Aprendizaje con Errores (LWE)}
    \begin{itemize}
        \item \textbf{Naturaleza:} Técnica asimétrica basada en problemas de retículos.
        \item \textbf{Clave Pública ($pk$):} Par $(\mathbf{A}, \mathbf{b})$ donde $\mathbf{b} = \mathbf{As} + \mathbf{e} \pmod q$.
    \end{itemize}
    
    \begin{block}{Ejemplo de Estructura}
        \[
        \underbrace{
        \begin{pmatrix} b_1 \\ b_2 \\ \vdots \\ b_m \end{pmatrix}
        }_{\mathbf{b} \in \mathbb{Z}_q^m}
        =
        \underbrace{
        \begin{pmatrix} 
        a_{11} & \dots & a_{1n} \\
        \vdots & \ddots & \vdots \\
        a_{m1} & \dots & a_{mn} 
        \end{pmatrix}
        }_{\mathbf{A} \in \mathbb{Z}_q^{m \times n}}
        \cdot
        \underbrace{
        \begin{pmatrix} s_1 \\ \vdots \\ s_n \end{pmatrix}
        }_{\mathbf{s} \in \mathbb{Z}_q^n}
        +
        \underbrace{
        \begin{pmatrix} e_1 \\ \vdots \\ e_m \end{pmatrix}
        }_{\mathbf{e} \in \chi^m} \pmod q
        \]
    \end{block}
    
    \begin{itemize}
        \item \textbf{Parámetros:} $n$ (seguridad), $m$ (muestras), $q$ (módulo), $\mathbf{e}$ (pequeño error).
    \end{itemize}
\end{frame}

\begin{frame}{Cifrado con LWE}
    Para cifrar un bit $m \in \{0, 1\}$ usando la clave pública $(\mathbf{A}, \mathbf{b})$:
    \begin{itemize}
        \item Seleccionar un vector aleatorio de muestreo $\mathbf{r}$.
        \item \textbf{Componentes del Criptograma $(c_1, c_2)$:}
        \[ \mathbf{c_1} = \mathbf{rA} \pmod q \]
        \[ c_2 = \mathbf{rb} + m\lfloor q/2 \rfloor \pmod q \]
        \item El mensaje se "esconde" desplazándolo a la mitad del módulo.
    \end{itemize}
\end{frame}

\begin{frame}{Descifrado con LWE}
    Para recuperar el mensaje $m$ usando la clave privada $\mathbf{s}$:
    \begin{itemize}
        \item \textbf{Operación de Descifrado:}
        \[ \mu = c_2 - \langle \mathbf{c_1}, \mathbf{s} \rangle \pmod q \]
        \item \textbf{Decisión:}
        \[ m = 
        \begin{cases} 
        0 & \text{si } \mu \text{ está más cerca de } 0 \\
        1 & \text{si } \mu \text{ está más cerca de } q/2 
        \end{cases} \]
        \item El descifrado es exitoso si el ruido acumulado es pequeño.
    \end{itemize}
\end{frame}

\begin{frame}{Aprendizaje con Errores en Anillos (RLWE)}
    El problema RLWE es una variante algebraica de LWE que traslada las operaciones de vectores y matrices a anillos de polinomios, donde la seguridad reside en la dificultad de distinguir muestras ruidosas de muestras aleatorias uniformes en el anillo.

    \begin{itemize}
        \item \textbf{Estructura:} Definido sobre el anillo $R_q = \mathbb{Z}_q[x] / \langle x^n + 1 \rangle$.
        \item \textbf{Ventajas:}
        \begin{itemize}
            \item Mayor eficiencia: Multiplicación vía \textbf{NTT}.
            \item Estructura compacta: Claves y textos cifrados más pequeños.
            \item Habilita operaciones paralelas nativas (\textbf{SIMD}).
        \end{itemize}
    \end{itemize}
\end{frame}

\begin{frame}{Esquemas FHE Modernos (Comparativa)}
    \begin{table}
        \centering
        \resizebox{\textwidth}{!}{
        \begin{tabular}{|l|l|l|l|}
            \hline
            \textbf{Característica} & \textbf{BGV / BFV} & \textbf{CKKS} & \textbf{TFHE} \\ \hline
            \textbf{Tipo de Dato} & Enteros exactos & Reales / Complejos & Bits / Enteros cortos \\ \hline
            \textbf{Gestión de Ruido} & Modulus/Scale switching & Reescalado & Bootstrapping rápido \\ \hline
            \textbf{Unidad de Cómputo} & Vectores (SIMD) & Vectores (SIMD) & Puertas Lógicas / LUT \\ \hline
            \textbf{Precisión} & Absoluta & Aproximada & Absoluta (Binaria) \\ \hline
        \end{tabular}
        }
    \end{table}
    \vfill
    \begin{itemize}
        \item \textbf{CKKS:} Seleccionado para este trabajo por su soporte nativo de aritmética aproximada, ideal para Redes Neuronales.
    \end{itemize}
\end{frame}
